% §3.5 Composition is the Bridge to Sets
%
% Cross-reference resolution status:
%
% RESOLVED:
%   \Cref{sec:and-under-noise}, \Cref{sec:composition-theorem},
%   \Cref{thm:composition}, \Cref{sec:miller-rabin}, \Cref{eq:mr-k-round} -- ch.3
%
% FORWARD REFS:
%   \Cref{ch:approximate-sets}      -- chapter 5
%   \Cref{ch:set-composition}       -- chapter 6
%   \Cref{ch:approximate-functions} -- chapter 9
%   \Cref{ch:not-bernoulli}         -- chapter 16
%   \Cref{sec:bib-notes-ch3}        -- chapter 3 §3.6
%
% Monograph register: three bridges from Bernoulli[Bool] to Part II/III.
% The intersection FPR product is the sparse/all-negative case; chapter 6
% carries the boundary-dependent general case (consistent with §3.1).

\section{Composition is the Bridge to Sets}
\label{sec:bridge-to-sets}

The Boolean algebra of $\Bernoulli{\BoolType}$ is the foundation for the
richer types of Parts~II and~III, and three concrete bridges carry it there.

\paragraph{Bloom filters are $m$-fold conjunctions.}
A Bloom filter reports membership by checking that all $m$ hash positions are
set, an $m$-fold AND of one-sided lookups. Each lookup has $\fnrate = 0$ (a
member sets every position) and per-hash false-positive rate $\fprate_i$. In
the all-negative regime of a non-member, the conjunction false-positives only
if every hash collides, so by the dual \eqref{eq:mr-k-round} of the composition
theorem the filter's false-positive rate is $\prod_{i=1}^{m} \fprate_i \le
p^m$, the classical Bloom formula, recovered as Boolean algebra over noisy
bits. The structure is identical to Miller-Rabin; only the source of the
per-stage rate differs, a design parameter here, a number-theoretic theorem
there. \Cref{ch:approximate-sets} develops the filter and its parameter
selection in full.

\paragraph{Set operations are Boolean operations on membership.}
For approximate sets $\tilde{A}, \tilde{B}$, membership $x \in \tilde{A}$ is a
$\Bernoulli{\BoolType}$ instance, and the set operations are pointwise Boolean
operations on these: intersection is AND, union is OR, complement is NOT. The
rate formulas of \Cref{sec:and-under-noise} and \Cref{sec:or-not} therefore
give, per element,
\[
  \fnrate_\cap = 1 - (1 - \fnrate_A)(1 - \fnrate_B),
  \quad
  \fprate_\cup = \fprate_A + \fprate_B - \fprate_A \fprate_B,
  \quad
  (\fprate_{A^c}, \fnrate_{A^c}) = (\fnrate_A, \fprate_A).
\]
The intersection's false-positive rate is the product $\fprate_A \fprate_B$ for
elements outside both sets (the bulk regime), and is boundary-dependent for
elements in exactly one set, exactly the prior-dependence of the AND in
\Cref{sec:and-under-noise}. \Cref{ch:set-composition} carries the general case,
including the rate-uniformity conditions an approximate set must satisfy for
these per-element rates to define a single $\Bernoulli{\Set{S}}$ instance,
conditions that chapter~16 (\Cref{ch:not-bernoulli}) shows can fail.

\paragraph{The Boolean case is a fiber of the map algebra.}
Replacing the Boolean codomain with an arbitrary type $Y$ lifts the same
composition law: an approximate map $\tilde{f} : X \to Y$ has a channel matrix
that grows from $2 \times 2$ to $|Y| \times |Y|$, but composition is still the
matrix product and the independence structure is unchanged. The Boolean algebra
of this chapter is the $Y = \BoolType$ fiber of the map algebra that
\Cref{ch:approximate-functions} develops. The composition law is the constant;
only the codomain varies.

The bibliographic notes (\Cref{sec:bib-notes-ch3}) follow. Chapter~4
formalizes the independence axioms these derivations assumed, and Part~II
builds the full algebra of approximate sets and maps on both chapters.
